\phantomsection
\section*{Глава 4. Оформление итогов работы}
\addcontentsline{toc}{section}{Глава 4. Оформление итогов работы}
\stepcounter{section}

В четвёртой главе описывается организация кодовой базы в системе контроля версий, тестирование и развёртывание приложения.

\subsection{Организация Git-репозитория}

Исходный код проекта размещён в публичном репозитории на платформе GitHub:

\begin{center}
\textbf{\url{https://github.com/Haskeri/kpweb}}
\end{center}

Развёрнутое веб-приложение доступно по адресу:

\begin{center}
\textbf{\url{https://exam.anufriev.pvpcult.ru/}}
\end{center}

Кодовая база проекта организована в репозитории Git. Корень репозитория содержит каталог \texttt{pcshop/} с Django-проектом, файл \texttt{requirements.txt} со списком зависимостей и файл \texttt{.gitignore}.

Репозиторий содержит полную историю разработки, сформированную в соответствии с этапами курсовой работы. История коммитов отражает последовательное выполнение разделов (таблица~\ref{tab:commits}).

\begin{table}[ht!]
\caption{История коммитов репозитория}
\label{tab:commits}
\small
\begin{tabularx}{\textwidth}{|p{0.22\textwidth}|p{0.18\textwidth}|>{\RaggedRight\arraybackslash}X|}
\hline
\textbf{Коммит} & \textbf{Дата} & \textbf{Содержание} \\ \hline
Раздел 1 & 25.04.2026 & Анализ предметной области, постановка задачи, структура LaTeX-проекта \\ \hline
Раздел 2 & 16.05.2026 & Проектирование архитектуры, схема БД, диаграммы компонентов \\ \hline
Раздел 3 & 06.06.2026 & Разработка веб-приложения, исходный код проекта pcshop (Django 4.2) \\ \hline
Раздел 4 & 13.06.2026 & Итоговый отчёт, развёртывание, оформление пояснительной записки \\ \hline
\end{tabularx}
\end{table}

Файл \texttt{.gitignore} исключает из репозитория:
\begin{itemize}
    \item виртуальное окружение \texttt{.venv/} и кэш Python \texttt{\_\_pycache\_\_/};
    \item базу данных \texttt{db.sqlite3} (содержит пользовательские данные);
    \item файл конфигурации \texttt{configs/.env} с реальными секретами;
    \item медиафайлы (загружаемые изображения комплектующих).
\end{itemize}

В репозитории хранится только файл \texttt{configs/.env.example} --- шаблон переменных окружения без реальных значений. Это обеспечивает безопасность: секреты и чувствительные настройки не попадают в историю коммитов.

Структура файла \texttt{requirements.txt} приведена в таблице~\ref{tab:requirements}.

\begin{table}[ht!]
\caption{Зависимости проекта (\texttt{requirements.txt})}
\label{tab:requirements}
\small
\begin{tabularx}{\textwidth}{|>{\RaggedRight\arraybackslash}p{0.35\textwidth}|>{\RaggedRight\arraybackslash}X|}
\hline
\textbf{Пакет} & \textbf{Назначение} \\ \hline
\texttt{Django==4.2.16} & Веб-фреймворк (LTS-версия до апреля 2026~г.) \\ \hline
\texttt{Pillow==10.3.0} & Обработка изображений (загрузка фото комплектующих) \\ \hline
\texttt{python-decouple==3.8} & Чтение конфигурации из файла \texttt{.env} \\ \hline
\texttt{psycopg2-binary==2.9.9} & Драйвер PostgreSQL (для продакшен-развёртывания) \\ \hline
\end{tabularx}
\end{table}

Установка зависимостей выполняется командой \texttt{pip install -r requirements.txt} в активированном виртуальном окружении Python 3.10+. Для инициализации базы данных и загрузки начальных данных предусмотрены команды:

\begin{lstlisting}[language=bash, caption={Команды инициализации проекта}]
python manage.py migrate
python manage.py loaddata fixtures/order_statuses.json
python manage.py loaddata fixtures/categories.json
python manage.py loaddata fixtures/components.json
python manage.py createsuperuser
\end{lstlisting}

\subsection{Тестирование и развёртывание приложения}

\textbf{Проверка системы.} Встроенная команда Django \texttt{python manage.py check} выполняет статическую проверку конфигурации, URL-маршрутов, моделей и установленных приложений. Команда \texttt{python manage.py check --deploy} дополнительно проверяет настройки безопасности (HTTPS, CSRF, SECRET\_KEY). Оба вида проверки не выявляют ошибок.

\textbf{Ручное тестирование.} Проведено функциональное тестирование всех ключевых сценариев:
\begin{itemize}
    \item регистрация нового клиента и вход по email;
    \item добавление комплектующих в корзину со страницы каталога и со страницы товара;
    \item создание сборки в конфигураторе с проверкой совместимости;
    \item оформление заказа с уменьшением остатков на складе;
    \item смена статуса заказа менеджером;
    \item экспорт и импорт каталога в формате CSV.
\end{itemize}

Все сценарии выполняются без ошибок на Python 3.9 при использовании SQLite в качестве СУБД.

\textbf{Развёртывание.} Для локального запуска используется встроенный сервер разработки:

\begin{lstlisting}[language=bash, caption={Запуск сервера разработки}]
cd pcshop
source .venv/bin/activate
python manage.py runserver
\end{lstlisting}

Приложение становится доступным по адресу \texttt{http://localhost:8000/}. Интерактивная административная панель доступна по адресу \texttt{http://localhost:8000/admin/}.

\textbf{Продакшен-развёртывание.} Помимо локального запуска, проект развёрнут на хостинге и доступен по адресу:

\begin{center}
\url{https://exam.anufriev.pvpcult.ru/}
\end{center}

На сервере используется стек: \textbf{Nginx} (обратный прокси) → \textbf{Gunicorn} (WSGI-сервер) → \textbf{Django}. Для корректной работы CSRF-защиты за HTTPS-прокси в настройках указан доверенный источник:

\begin{lstlisting}[language=bash, caption={Ключевые параметры .env на сервере}]
DEBUG=False
ALLOWED_HOSTS=exam.anufriev.pvpcult.ru
CSRF_TRUSTED_ORIGINS=https://exam.anufriev.pvpcult.ru
SECRET_KEY=<длинный случайный ключ>
\end{lstlisting}

Переход на продакшен-окружение предполагает замену SQLite на PostgreSQL через переменные окружения в \texttt{.env}:

\begin{lstlisting}[language=bash, caption={Пример файла .env для PostgreSQL}]
SECRET_KEY=your-long-random-key
DEBUG=False
ALLOWED_HOSTS=yourdomain.com
DB_ENGINE=django.db.backends.postgresql
DB_NAME=pcshop
DB_USER=pcshop_user
DB_PASSWORD=secure_password
DB_HOST=localhost
DB_PORT=5432
\end{lstlisting}

Такой подход обеспечивает принцип «12-factor app»: конфигурация полностью отделена от кода и не попадает в репозиторий.

\subsection*{Выводы по главе 4}
\addcontentsline{toc}{subsection}{Выводы по главе 4}

В четвёртой главе описана организация кодовой базы в системе Git с корректным \texttt{.gitignore}, разделением конфигурации через \texttt{.env} и фикстурами для первоначального заполнения базы данных. Исходный код размещён в публичном репозитории по адресу \url{https://github.com/Haskeri/kpweb}. Проведено ручное функциональное тестирование всех ключевых пользовательских сценариев. Приложение развёрнуто на хостинге и доступно по адресу \url{https://exam.anufriev.pvpcult.ru/}. Описаны команды запуска в режиме разработки и инструкция по переходу на продакшен-СУБД PostgreSQL без изменения кода приложения.
